\SourceLeaf{407}%
until I have gathered all the materials for these subjects. That will take time, partly because there are so many topics and partly because the difficulties keep multiplying and new ones keep appearing. My thinking on the matter will now again be interrupted by some lecturing.
\par\Needspace{8\baselineskip}
\begin{center}\textsc{Gauss} to \textsc{Schumacher}. Göttingen, 21 November 1825.\end{center}
\indent\mbox{. . . . . . .} For some time I have been working again on part of the general investigations of curved surfaces that are to form the basis of my projected work on Higher Geodesy. It is as rich a subject as it is difficult, and at present it keeps me from getting to any other work. Unfortunately I find that I shall have to go \textls{very far back}, since even the familiar material must be developed in a different form adapted to the new investigations. One must follow the tree down to every one of its root fibres, and some of this costs me weeks of strenuous thought. Much of it even belongs to the Geometria situs, a field still almost untouched. My constant desire to bring my work to such completion, ut nihil amplius desiderari possit, naturally makes it extraordinarily difficult for me, as does the need often to break away from it for unrelated matters. When my head is full of it, you can scarcely imagine how taxing it sometimes is for me in the morning, after a sleepless night, which I unfortunately often have now, to turn with a fresh mind to what I have to present to my students, and afterward immediately to return with renewed energy to my meditations. From time to time, however, I am rewarded by a happy new discovery. For example, I have now brought to an exceedingly elegant form my generalization of \textsc{Legendre}'s theorem, according to which on the sphere the sides are proxime proportional to the sines of the angles
